\chapter{Open form-identification bridges}
\label{chap:open-canonical-bridges}

Closed quadratic structure is only the analytic base.  The following form-level bridges are kept
as named theorems so later work cannot smuggle them into the definition of the selected canonical
constant.  The constant-level bridges used by every route are collected separately in
\cref{chap:shared-analytic-bridges}.

\begin{obstruction}[Raw energy is not a property of an $L^2$ class]
  \label{obs:dirac-raw-energy}
  Let $\mu=\delta_0$ on $\R$, $f(x)=0$, and $g(x)=x$.  Then $f=g$ in $L^2(\mu)$, but their raw
  gradient energies are respectively $0$ and $1$.  Hence unconditional agreement of raw test
  energies cannot descend to $L^2$ equivalence classes for singular measures.
\end{obstruction}
\begin{proof}
  Both functions take value zero at the support point, while $f'=0$ and $g'=1$ everywhere.
\end{proof}

\section{Form properties and literature normalization}

\begin{theorem}[Dirichlet property of the smooth relaxation]
  \label{thm:smooth-gradient-relaxation-dirichlet}
  \uses{thm:smooth-gradient-relaxation-closed}
  For every finite-dimensional Borel probability $\mu$, the ambient smooth-gradient relaxation
  is Markovian: every normal contraction of an $L^2(\mu)$ class decreases its relaxed energy.
  Together with closed quadratic structure, it therefore defines a Dirichlet energy.
\end{theorem}
\begin{proof}
  Open.  The reusable contraction-transfer theorem already isolates the required input; the
  missing step is a concrete smooth approximation argument uniform in the relaxed energy.
\end{proof}

\begin{theorem}[Identification with Cheeger energy]
  \label{thm:smooth-gradient-relaxation-two-cheeger}
  \notready
  \uses{def:selected-canonical-frontend}
  On the affine support, let $\Ch_\mu$ denote the standard metric Cheeger energy normalized as
  one half of the lower-semicontinuous relaxation of the squared local slope.  After transporting
  it to ambient real $L^2(\mu)$, prove
  \[
    \mathcal E^{\mathrm{can}}_\mu=2\Ch_\mu.
  \]
\end{theorem}

The factor $2$ is part of the intended normalization, not an established declaration.  This node
remains \texttt{notready} until a local Cheeger-energy API and its affine-support convention are
frozen; it does not presuppose that Mathlib will supply that API.

\section{Core and coordinate comparisons}

\begin{theorem}[Raw tests agree with the closed core]
  \label{thm:smooth-gradient-raw-core-agreement}
  \notready
  \uses{obs:dirac-raw-energy,def:smooth-gradient-preenergy,def:smooth-gradient-relaxation}
  Under hypotheses that exclude representative-dependent raw energy, formulate and prove the
  exact \texttt{AgreesOnTests} and form-core statement relating smooth tests to the relaxed
  energy.
\end{theorem}

The Dirac obstruction shows why the missing hypotheses are mathematical data.  Neither the
Dirichlet property nor a Cheeger identification is declared as a dependency until an actual proof
uses it.  Raw/core agreement is an optional, strictly stronger testwise identification.  It is
not a premise of the legacy--relaxed constant bridge.

\begin{theorem}[Ambient and intrinsic relaxations agree]
  \label{thm:ambient-intrinsic-relaxation-equality}
  \notready
  \uses{def:smooth-gradient-relaxation,def:intrinsic-affine-energy,thm:intrinsic-l2-equivalence}
  For every finite-dimensional Borel probability $\mu$, the intrinsic energy is already pulled
  back to ambient $L^2(\mu)$, and on that common carrier
  \[
    \mathcal E^{\mathrm{amb}}_\mu(f)
      =\mathcal E^{\mathrm{intr}}_\mu(f)
    \qquad\text{for every }f\in L^2(\mu).
  \]
  Consequently the ambient and selected intrinsic-affine optimal constants are equal.  The
  equality is independent of the affine origin used to write intrinsic coordinates.
\end{theorem}
\begin{proof}
  Open.  Restriction of an ambient smooth test to the affine support projects its gradient onto
  the tangent space and cannot increase energy.  Conversely, composing an intrinsic smooth test
  with the orthogonal projection onto the affine support gives a global ambient smooth extension
  with the same energy.  Apply both constructions to approximating sequences and take the two
  relaxations.
\end{proof}

The preferred implementation order is the exact ambient/intrinsic equality, the direct
legacy--relaxed bridge of \cref{chap:shared-analytic-bridges}, and then the literature test-density
bridge.  Dirichlet and Cheeger identifications follow as form enhancements.  A raw/core theorem
with explicit singular-measure hypotheses remains useful, but it is not on this main path.
